% !TeX root = ../../main.tex
% =====================================================================
%  Formula Student: el juego y sus reglas
%  Responsable:
%  Última revisión:
% =====================================================================
\chapter{Formula Student: el juego y sus reglas}
\label{ch:formula_student}

\begin{objetivos}
  \item Enumerar los eventos de una competición y cuántos puntos vale cada uno.
  \item Explicar por qué la fiabilidad da más puntos que las prestaciones.
  \item Localizar en el reglamento el artículo que afecta a tu subsistema y saber
        qué hacer cuando es ambiguo.
  \item Reconocer los documentos de temporada y las consecuencias de entregarlos tarde.
  \item Ordenar por prioridad el trabajo del equipo en función de dónde están los puntos.
\end{objetivos}

\fichasesion{90 minutos}{Ninguno. Es la primera sesión del curso}%
{Reglamento vigente descargado y calendario de la temporada}

Formula Student no es una carrera. Es una competición de \emph{ingeniería} en la
que hay, además, unas pruebas en pista. Esta distinción no es un matiz retórico:
determina cómo se reparten los puntos, y por tanto cómo debe repartir el equipo
su tiempo. Un tercio de la puntuación se decide hablando con jueces, sin
encender el coche.

Este capítulo describe cómo se reparten los puntos y qué documentos hay que
entregar durante la temporada. El resto del libro se apoya en ello: las
decisiones de diseño de los capítulos siguientes se justifican, en última
instancia, contra este reparto.

% ---------------------------------------------------------------------
\section{La competición y sus eventos}
\label{sec:formula_student:01}
\index{eventos}

Una competición se divide en \textbf{eventos estáticos} y \textbf{eventos
dinámicos}, con una barrera entre medias: la \textbf{revisión técnica}
(\emph{scrutineering}).

\subsection{Eventos estáticos}

Se disputan con el coche parado, delante de un panel de jueces que son
ingenieros del sector.

\begin{description}[leftmargin=0pt,labelsep=0pt,style=nextline]
  \item[\textit{Business Plan Presentation}] El equipo presenta un plan de
    negocio para el vehículo como si fuera un producto real. No se juzga la
    ingeniería, sino la coherencia del modelo de negocio y la calidad de la
    presentación.
  \item[\textit{Cost and Manufacturing}] Se entrega un informe con el coste
    completo del coche calculado según unas tablas comunes a todos los equipos, y
    se defiende ante los jueces. Se evalúa la precisión del informe, el
    conocimiento real de los procesos de fabricación y, en el \textit{Real Case
    Scenario}, la capacidad de razonar sobre coste bajo un supuesto que dan en el
    momento. Se trata en detalle en el capítulo \ref{ch:evento_costes}.
  \item[\textit{Engineering Design}] El evento más importante de los tres y el
    que más puntos reparte. Los jueces preguntan \emph{por qué} está diseñado así
    cada subsistema. No buscan el coche más sofisticado: buscan decisiones
    justificadas con datos y coherentes con los objetivos que el propio equipo se
    ha fijado. Se prepara con el capítulo \ref{ch:design_event}.
\end{description}

\subsection{Revisión técnica}
\index{revisión técnica}\index{scrutineering}

No reparte puntos. Es una \emph{puerta}: hasta que el coche no la pasa, no puede
salir a pista. Incluye, según el tipo de vehículo y la competición, la inspección
mecánica y eléctrica, la prueba de vuelco, el ensayo de ruido o de estanqueidad,
la prueba de salida del piloto (\emph{egress}) y el ensayo de frenada.

\begin{clave}
Superar la revisión técnica vale, en la práctica, los \textbf{675 puntos
dinámicos}: sin la pegatina, todos ellos son cero. Es la actividad con mayor
rentabilidad por hora invertida de toda la temporada.
\end{clave}

\subsection{Eventos dinámicos}

\begin{description}[leftmargin=0pt,labelsep=0pt,style=nextline]
  \item[\textit{Acceleration}] Recta de \SI{75}{\meter} desde parado. Mide
    tracción y relación de transmisión; el aerodinámico aquí solo estorba.
  \item[\textit{Skidpad}] Un ocho de radio constante. Mide agarre lateral
    estacionario y equilibrio del coche.
  \item[\textit{Autocross}] Una vuelta rápida a un circuito estrecho y revirado.
    Además de puntuar, fija el orden de salida de la resistencia.
  \item[\textit{Endurance}] \SI{22}{\kilo\meter} con un cambio de piloto a
    mitad. Es el evento decisivo: mide fiabilidad y, de paso, ritmo sostenido.
  \item[\textit{Efficiency}] No se corre: se calcula con el consumo y el tiempo
    de la resistencia. \emph{Solo puntúa si se termina la resistencia.}
\end{description}

% ---------------------------------------------------------------------
\section{El reparto de puntos}
\label{sec:formula_student:02}
\index{puntuación}

\begin{table}[H]
\centering
\caption{Reparto habitual de la puntuación. Verificar siempre contra el
reglamento vigente: los valores cambian entre categorías y entre temporadas.}
\label{tab:reparto_puntos}
\begin{tabular}{@{}llS[table-format=3.0]@{}}
\toprule
\textbf{Evento} & \textbf{Tipo} & {\textbf{Puntos}} \\
\midrule
Business Plan Presentation & Estático & 75 \\
Cost and Manufacturing     & Estático & 100 \\
Engineering Design         & Estático & 150 \\
\midrule
\multicolumn{2}{@{}l}{\textbf{Subtotal estáticos}} & \bfseries 325 \\
\midrule
Acceleration & Dinámico & 75 \\
Skidpad      & Dinámico & 75 \\
Autocross    & Dinámico & 100 \\
Endurance    & Dinámico & 325 \\
Efficiency   & Dinámico & 100 \\
\midrule
\multicolumn{2}{@{}l}{\textbf{Subtotal dinámicos}} & \bfseries 675 \\
\midrule
\multicolumn{2}{@{}l}{\textbf{Total}} & \bfseries 1000 \\
\bottomrule
\end{tabular}
\end{table}

Conviene mirar esta tabla con atención, porque contiene casi toda la estrategia
del equipo:

\begin{itemize}
  \item \textbf{Los estáticos son 325 puntos que no dependen de que el coche
        funcione.} Se pueden preparar con el coche a medio montar, y de hecho es
        lo que hacen los equipos que puntúan bien en ellos.
  \item \textbf{\textit{Endurance} y \textit{Efficiency} suman 425 puntos}, más
        que todos los estáticos juntos. Y están acoplados: sin terminar la
        resistencia, ambos son cero.
  \item \textbf{\textit{Acceleration} y \textit{Skidpad} suman 150 puntos.} Son
        los eventos donde más se nota el rendimiento puro y, sin embargo, los que
        menos pesan. Es un error clásico diseñar el coche para ellos.
\end{itemize}

Además, la puntuación de cada evento dinámico no es lineal: se reparte por
comparación con el mejor equipo y con un tiempo máximo, de modo que
\textbf{terminar despacio puntúa mucho más que no terminar}. La diferencia entre
el primero y el último clasificado \emph{de los que terminan} es siempre menor
que la diferencia entre terminar y no terminar.

% ---------------------------------------------------------------------
\section{Cómo se lee el reglamento}
\label{sec:formula_student:03}
\index{reglamento}

El reglamento no es un anexo administrativo: es el primer documento de
requisitos del proyecto, y llega escrito antes de que empieces a diseñar. Buena
parte de las decisiones de arquitectura del coche están ya tomadas en él.

\subsection{Estructura}

Las normas están agrupadas por secciones identificadas con letras, y cada
artículo se cita con la letra de su sección y una numeración jerárquica (por
ejemplo, \texttt{T 2.1.1}). Las secciones habituales son:

\begin{center}
\begin{tabular}{@{}ll@{}}
\toprule
\textbf{Sección} & \textbf{Contenido} \\
\midrule
A  & Normas administrativas: inscripción, plazos, documentos, penalizaciones \\
T  & Normas técnicas generales, aplicables a todos los vehículos \\
C  & Específicas de vehículos de combustión \\
EV & Específicas de vehículos eléctricos \\
DV & Específicas de vehículos autónomos \\
IN & Inspecciones y revisión técnica \\
S  & Eventos estáticos \\
D  & Eventos dinámicos \\
\bottomrule
\end{tabular}
\end{center}

\begin{regla}[cómo citar]
En este libro, y en cualquier documento del equipo, un requisito normativo se
cita \emph{siempre} con su artículo y la versión del reglamento:
«\textit{la distancia mínima entre el casco y la línea que une los arcos es de
\SI{50}{\milli\meter} (T 1.4.2, FS Rules 2025 v1.1)}». Sin artículo no es un
requisito: es una opinión.
\end{regla}

\subsection{Método de trabajo}

\begin{enumerate}
  \item \textbf{Léelo entero una vez, al principio de la temporada}, incluidas
        las secciones que no son de tu subsistema. Muchos requisitos afectan a
        varios a la vez.
  \item \textbf{Extrae los artículos de tu subsistema a una tabla de
        trazabilidad} con cuatro columnas: artículo, qué exige, cómo lo
        cumplimos, dónde se demuestra. Esa tabla es simultáneamente la
        preparación de la revisión técnica y media memoria de diseño.
  \item \textbf{Distingue lo obligatorio de lo recomendado.} El reglamento usa un
        lenguaje deliberadamente preciso; una frase con «\textit{must}» o
        «\textit{shall}» no admite interpretación.
  \item \textbf{Si algo es ambiguo, pregunta por el cauce oficial.} Las
        competiciones tienen un sistema de consultas cuyas respuestas son
        vinculantes. Una respuesta por escrito vale en la revisión técnica; la
        interpretación de un compañero, no.
  \item \textbf{No copies números de la temporada anterior.} El reglamento cambia
        todos los años, y las cotas, espesores y materiales admisibles están
        entre lo que más cambia.
\end{enumerate}

\begin{ojo}
\begin{itemize}
  \item Diseñar primero y buscar el artículo después, para «comprobar que
        cumple». El orden correcto es el inverso.
  \item Leer solo la sección propia. Los requisitos de accesibilidad,
        protecciones y separaciones son transversales y se incumplen en las
        \emph{fronteras} entre subsistemas.
  \item Dar por bueno el diseño del coche anterior. Que pasara la revisión hace
        dos temporadas no significa que la pase con el reglamento de esta.
\end{itemize}
\end{ojo}

% ---------------------------------------------------------------------
\section{Documentos y plazos de la temporada}
\label{sec:formula_student:04}
\index{documentos de temporada}

Una parte importante de la competición se juega meses antes de llegar al
circuito, entregando documentos. Casi todos tienen penalización por retraso, y
algunos son excluyentes: sin ellos no se compite.

\begin{table}[H]
\centering
\caption{Documentos habituales de temporada. Los plazos concretos y la lista
exacta dependen de la competición: rellenar al inicio de cada temporada.}
\label{tab:documentos_temporada}
\begin{tabularx}{\textwidth}{@{}lXll@{}}
\toprule
\textbf{Documento} & \textbf{Qué es} & \textbf{Responsable} & \textbf{Fecha} \\
\midrule
Inscripción y cuestionario & Prueba de conocimiento del reglamento & & \\
SES & Justificación de equivalencia estructural del chasis & & \\
IAD & Diseño y ensayo del atenuador de impactos & & \\
ESF & Formulario del sistema eléctrico (vehículos eléctricos) & & \\
Cost Report y BOM & Coste completo del vehículo & & \\
Engineering Design Report & Memoria de diseño & & \\
Design Spec Sheet & Ficha resumen de características & & \\
Business Plan & Documentación del evento de negocio & & \\
Vídeo de estado del vehículo & Prueba de que el coche rueda & & \\
\bottomrule
\end{tabularx}
\end{table}

\pendiente{Rellenar la tabla con los documentos, responsables y fechas reales de
la temporada en curso, y colgarla donde se vea todos los días.}

\begin{clave}
Los plazos de documentación son la única parte de la competición donde se pierden
puntos \emph{sin haber hecho nada mal técnicamente}. Cuestan cero horas de taller
y solo exigen organización. Un equipo que pierde puntos aquí está regalando la
parte más barata de la competición.
\end{clave}

% ---------------------------------------------------------------------
\section{De dónde salen realmente los puntos}
\label{sec:formula_student:05}

Si se cruza el reparto de puntos con lo que ocurre de verdad en una competición,
sale una lista de prioridades bastante distinta de la intuitiva. Un porcentaje
muy alto de los equipos inscritos no termina la resistencia, y una parte
importante ni siquiera llega a pasar la revisión técnica a tiempo de disputar
todos los eventos.

De ahí el orden de prioridad con el que debería trabajar cualquier equipo, y muy
especialmente uno joven:

\begin{enumerate}
  \item \textbf{Pasar la revisión técnica.} Diseñar desde el primer día con la
        lista de comprobación de la inspección delante, y hacer una inspección
        interna simulada semanas antes de viajar.
  \item \textbf{Terminar la resistencia.} Kilómetros de rodaje antes de la
        competición, componentes con margen, repuestos y procedimientos de
        montaje. La fiabilidad no se añade al final: se diseña.
  \item \textbf{Preparar bien los estáticos.} 325 puntos que solo cuestan
        organización, documentación y ensayo de la defensa.
  \item \textbf{Y después, ir rápido.} Cuando lo anterior está resuelto, las
        prestaciones deciden posiciones. Antes, no deciden nada.
\end{enumerate}

\begin{clave}
El coche no tiene que ser rápido: tiene que \textbf{terminar}. Un coche lento que
completa los cinco eventos dinámicos y defiende bien sus estáticos puntúa por
encima de un coche brillante que rompe en la resistencia.
\end{clave}

Esta jerarquía tiene una consecuencia incómoda para un equipo de ingenieros: casi
siempre, la decisión correcta es \emph{la más simple que cumple el objetivo}. La
sofisticación solo se justifica cuando el equipo ya es capaz de fabricarla,
validarla y repararla en el paddock. Los jueces del evento de diseño valoran
exactamente eso, y penalizan la complejidad que el equipo no puede sostener.

% ---------------------------------------------------------------------
%  Cierre del capítulo
% ---------------------------------------------------------------------

\begin{caso}
\redactar{Rellenar con nuestro historial: puntuación obtenida por evento en las
últimas participaciones, en qué eventos no pudimos puntuar y por qué, y cuánto
tiempo tardamos en pasar la revisión técnica. Es el punto de partida honesto
para fijar los objetivos de esta temporada.}
\end{caso}

\begin{ojo}
\begin{itemize}
  \item Medir el éxito del equipo por el tiempo en \textit{acceleration}: 75 de
        1000 puntos.
  \item Empezar a preparar los estáticos cuando el coche ya rueda. Para entonces
        no queda tiempo, y son 325 puntos.
  \item Tratar la revisión técnica como un trámite del final en lugar de como un
        requisito de diseño.
  \item Fijar los objetivos de la temporada sin mirar la puntuación de la
        anterior.
\end{itemize}
\end{ojo}

\begin{entregable}
\begin{enumerate}
  \item La tabla \ref{tab:documentos_temporada} rellena con los plazos reales de
        esta temporada y un responsable por documento.
  \item Una tabla de trazabilidad de reglamento para tu subsistema: artículo,
        qué exige, cómo lo cumplimos, dónde se demuestra.
  \item Los objetivos de puntuación del equipo para esta temporada, evento a
        evento, con su justificación.
\end{enumerate}
\end{entregable}

\begin{juez}
\begin{enumerate}
  \item ¿Qué objetivos de puntuación se fijó el equipo esta temporada y por qué esos?
  \item ¿Qué evento habéis priorizado en el diseño del coche, y qué habéis
        sacrificado a cambio?
  \item ¿Qué artículo del reglamento ha condicionado más el diseño de tu
        subsistema? Cítalo.
  \item ¿Cuántos kilómetros ha rodado el coche antes de venir aquí?
  \item ¿Qué os impidió puntuar en algún evento la temporada pasada y qué habéis
        cambiado para que no vuelva a pasar?
\end{enumerate}
\end{juez}
